\section{Conclusion and Future Work}
\label{sec:conclusion}

We studied automatic scoring of flea beetle feeding damage on young oilseed rape from field
photographs with very few reliable labels. The three families of approaches behave very
differently.

Classical computer vision is reliable for the geometric steps, i.e.\ finding the frame and the
plants and measuring their visible area, but not for the damage itself: holes, soil and brown
pitting cannot be separated by colour, and eaten leaf edges do not show up in a pixel count. We
therefore see it mainly as preprocessing and as a source of interpretable side features such as
pitting area. Whole-image regression with a frozen DINOv3 fails (test $\rho = -0.11$), because
downsampling a field photograph to $224\times224$ erases the seedlings. Plant-level MIL works
in-domain, with a test MAE of 2.4--2.6 percentage points and $\rho$ of 0.76--0.81. Fixed pooling
rules were at least as good as learned attention. Area weighting won on validation and uniform
pooling on test; given how the raters score, uniform pooling may be closer to the label
definition. Adding a pairwise ranking loss gave a small improvement for all three seeds.
On the Gro\ss-Gerau images where the two raters disagreed, the model agrees with each rater
better than they agree with each other. On new field trials, however, the model does not
generalize ($\rho$ 0.03 and 0.26), mainly because the frame is not found and plants outside the
frame are scored, and partly because the learned features transfer poorly to other fields.

\paragraph{Limitations.} All reliable labels come from one field on one day, so our in-domain
numbers describe a single recording condition. The labels are subjective averages that mix holes
and pitting, which limits the achievable accuracy and makes the choice of pooling rule ambiguous.
Validation and test sets are small (66 and 73 images), so differences below about 0.1 in $\rho$
between methods are uncertain. Only 40 images were annotated for the hole/pitting detector, and
only 7 of the 217 genotypes have more than one plot in the calibration set, so we did not attempt
a genotype resistance ranking.

\paragraph{Recommendations for future data recording.} Most of the generalization failure could be
avoided at recording time:
(i) use the same frame and hold the camera so that the frame fills the photo, at a fixed height,
at every site;
(ii) give the frame a marker that is easy to detect, e.g.\ coloured corners or an ArUco marker next
to the QR code;
(iii) score a small calibration set (50--100 images) with two raters at every site and date, so that
each new domain has some reliable labels;
(iv) record, per image, whether single plants or the image as a whole were scored, and score holes
and pitting separately;
(v) annotate 150--300 images per damage class with boxes or masks to make a lesion detector usable.

\paragraph{Future work.} The next steps are, in order of expected benefit: a more robust frame and
plant detector (e.g.\ SAM-based segmentation of plants, as suggested in the project brief) and
re-running the out-of-distribution evaluation; few-shot calibration of the head with a few
labelled images per new site; domain adaptation by self-supervised pre-training on the 8{,}946
unlabelled images; weak labels from a vision-language model; and combining the ranking loss with
uniform pooling. A learned leaf and lesion segmentation would also allow the plant counts, leaf
counts and hole/pitting statistics requested by the project partners.
